\section{Proofs for Section~\ref{sec:ext}}\label{app:proofs-ext}

Notation is as in Section~\ref{sec:ext}. Throughout, the minimum over an empty set of modes is $+\infty$, $(x)_+=\max\{x,0\}$, and $\hat w_j\ge0$ and $\mathcal C_{x,j}\ge0$ by the assumptions stated after \eqref{eq:ext-wedges}. Two facts are used repeatedly.

\begin{enumerate}
\item[(E1)] (Ordering.) Let $\theta<\theta'$ and suppose that $\mathcal M_j$ does not depend on $\theta$ and that
\[
\inf_{x\in\mathcal M^{\rm AI}\cap\mathcal M_j}\bigl[\tilde c_x(\theta')-\tilde c_x(\theta)\bigr]\;\ge\;\sup_{x\in\mathcal M^{0}\cap\mathcal M_j}\bigl[\tilde c_x(\theta')-\tilde c_x(\theta)\bigr].
\]
Then $\mathcal B_j(\theta')\le\mathcal B_j(\theta)$. Indeed, write $\underline a$ for the infimum and $\bar b$ for the supremum. For every AI mode $\tilde c_x(\theta')\ge\tilde c_x(\theta)+\underline a$, so $\min_{\rm AI}\tilde c(\theta')\ge\min_{\rm AI}\tilde c(\theta)+\underline a$. For every non-AI mode $\tilde c_x(\theta')\le\tilde c_x(\theta)+\bar b$, so $\min_0\tilde c(\theta')\le\min_0\tilde c(\theta)+\bar b$. Subtracting, $\mathcal B_j(\theta')\le\mathcal B_j(\theta)+\bar b-\underline a\le\mathcal B_j(\theta)$. If instead a change removes AI modes from $\mathcal M_j$ or adds cost to them only, the AI minimum weakly rises and the non-AI minimum is unchanged (mode H is always available), so again $\mathcal B_j$ weakly falls.
\item[(E2)] (Danskin.) If each $\tilde c_x$ is differentiable in $\theta$ and the minimizers $x^0_j$, $x^1_j$ are unique at $\theta_0$, then, because the mode sets are finite, each minimum equals the minimizer's cost on a neighborhood of $\theta_0$, and $\partial\mathcal B_j/\partial\theta=\partial\tilde c_{x^0_j}/\partial\theta-\partial\tilde c_{x^1_j}/\partial\theta$.
\end{enumerate}

\paragraph{Definition~\ref{def:ext-effcost} (derivation).} Per unit in mode $x$, the owner bears the after-tax operating cost $(1-\tau)[w_jt_x+w^L_jt^L_x+c^{\rm ext}_x+(1+\tau_A)c^{\rm mach}_x]$. Capital spending, net of the present value of its deductions, costs $(1-\psi\tau)(1+\tau_A)c^K_x$. The penalty adds $\omega_x$ and altruism subtracts $\vartheta\bar t_x$, both in after-tax dollars. Lost contribution from customers is taxable, so it costs $(1-\tau)$ times its pre-tax amount. Dividing the sum by $1-\tau$ gives \eqref{eq:ext-effcost}, with $\hat w=w^0(1+\tau_L)-\vartheta/(1-\tau)$ and $\Upsilon=(1-\psi\tau)/(1-\tau)$. In mode A the owner also chooses the review share; after division by $1-\tau$ the terms that depend on it are $\hat w_j\xi t^R_j+e_jL_j[1-\xi(1+\ln(1/\xi))]$, the quality cost of Proposition~\ref{prop:tasks-review} at the wage $\hat w_j$, and the reliability and floor constraints on $\xi$ do not involve the wage. Its minimum $\hat c^Q_j$ is therefore given by \eqref{eq:tasks-xi} with $\hat w_j$ in place of $w_j$. It is the minimum over $\xi$ of functions affine in $\hat w_j$, with a unique minimizer because $c^Q_j$ is strictly convex in $\xi$, so by Danskin's theorem $\partial\hat c^Q_j/\partial\hat w_j=\xi^*_jt^R_j$. The pre-tax amount of the customer term is derived in the proof of Proposition~\ref{prop:ext-customer}. Finally, $\Upsilon-1=\tau(1-\psi)/(1-\tau)\ge0$, with equality if and only if $\psi=1$ or $\tau=0$. With straight-line deductions $1/T_R$ a year for $T_R$ years, $\psi=\int_0^{T_R}e^{-rt}T_R^{-1}dt=(1-e^{-rT_R})/(rT_R)$. The capital recovery factor satisfies $a(r,T)\int_0^Te^{-rt}dt=1$, and $a(r,T)\to1/T$ as $r\to0$. \qed

\paragraph{Lemma~\ref{lem:ext-nested}.} (a) Every AI mode's machine and capital terms are nonnegative, with positive coefficients $1+\tau_A$ and $\Upsilon$. Setting them to zero therefore weakly lowers every AI cost and leaves the non-AI costs unchanged, so $\mathcal B_j\le\mathcal B^0_j$ and $\mathbb A\subseteq\mathcal J^{\rm pre}$. The inclusion $\mathcal J^{\rm pre}\subseteq\mathcal J^{\rm per}$ holds by definition, and $\mathcal J^{\rm per}\subseteq\mathcal J^{\rm fea}$ holds because $\mathcal M_j\subseteq\mathcal M^F_j$. The spend shares inherit the inclusions because the weights are nonnegative, and \eqref{eq:ext-waterfall} telescopes. (b) $\mathcal B^0_j$ is computed with the machine and capital terms of the AI modes set to zero. The parameters $\pi$, $F$, $K$ and $\tau_A$ enter $\tilde c$ only through these terms, since $\tau_A$ is zero for the non-AI modes, so $\mathcal J^{\rm pre}$ does not depend on them. By (E1), $\mathcal B_j$ is nonincreasing in each of these costs, so $\mathbb A$ grows as they fall. Since $\mathcal B_j$ is continuous in them (a difference of minima of finitely many continuous functions), $\mathcal B_j\to\mathcal B^0_j$ as they tend to zero. Every $j$ with $\mathcal B^0_j>0$ therefore eventually lies in $\mathbb A$, so $\mathbb A\uparrow\mathcal J^{\rm pre}$. The last claim follows from (a), and from the fact that energy costs enter only $c^{\rm mach}$. \qed

\paragraph{Proposition~\ref{prop:ext-statics}.} (a) Each parameter listed raises the cost of AI modes only, or removes AI modes, and leaves non-AI costs unchanged. For $\tau_A$, $\pi$, $F_A$ and $K_A$ the increase is $(c^{\rm mach}+\Upsilon c^K)\,d\tau_A\ge0$, $(1+\tau_A)n\,d\pi\ge0$, and so on. For the floor, the quality cost $c^Q_j$ of Proposition~\ref{prop:tasks-review} is convex in the review share and minimized at the share chosen without the floor, so it is nondecreasing in a binding floor. For a penalty it is $d\omega/(1-\tau)\ge0$, and for the demand terms $\hat X(S/v)\,d\delta\ge0$. Exclusions, class 3 and licensing remove modes or add sign-off cost. A higher class raises the floor. By (E1), $\mathcal B_j$ is nonincreasing in each of these, so membership of $\mathbb A$ is monotone. $\Sigma^{\rm eco}$ is a sum of nonnegative weights over $\mathbb A$, and $\mathcal V^{\rm eco}$ a sum of $v_j(\mathcal B_j)_+$, so both are monotone. A higher $\psi$ lowers $\Upsilon$ and therefore lowers AI costs only, because capital charges are attached to AI modes. Higher readiness lowers the raw error rate $e_j(\iota)$, and with it the quality cost (Proposition~\ref{prop:tasks-review}(c)), and it can only enlarge $\mathcal M_j$; the same argument applied in reverse gives monotonicity.

(b) By (E2), each derivative is the difference of the minimizers' derivatives. From \eqref{eq:ext-effcost}, $\partial\tilde c_x/\partial\tau_L=w^0t_x+w^{L0}t^L_x$ and $\partial\tilde c_x/\partial\vartheta=-\bar t_x/(1-\tau)$, which gives \eqref{eq:ext-dBdL}; for mode A these follow from $\partial\hat c^Q_j/\partial\hat w_j=t_{\mathrm A}$ and the chain rule. For $\tau$, $\partial\hat w/\partial\tau=-\vartheta/(1-\tau)^2$, $\partial\Upsilon/\partial\tau=(1-\psi)/(1-\tau)^2$, $\partial[\omega/(1-\tau)]/\partial\tau=\omega/(1-\tau)^2$ and, holding $X$ and $k$ fixed, $\partial\hat X/\partial\tau=[X-(1-\psi)k]/(1-\tau)^2$. Non-AI modes carry no capital charge, penalty or customer term, so
\[
\frac{\partial\mathcal B_j}{\partial\tau}=-\frac{\vartheta\bar t_{x^0}}{(1-\tau)^2}-\Bigl[-\frac{\vartheta\bar t_{x^1}}{(1-\tau)^2}+\frac{(1+\tau_A)(1-\psi)c^K_{x^1}+\omega_{x^1}+\delta_{x^1}(S/v)(X-(1-\psi)k)}{(1-\tau)^2}\Bigr],
\]
which is \eqref{eq:ext-dBdtau}. Every term in the numerator is nonnegative when $\bar t_{x^0}\ge\bar t_{x^1}$, $\omega\ge0$ and $X\ge(1-\psi)k$, and all of them vanish when $\psi=1$ and $\omega=\vartheta=\delta=0$. The sign statements about $\tau_L$ and $\vartheta$ read off \eqref{eq:ext-dBdL}. For a global statement: the change in each $\tilde c_x$ between two values of $\tau_L$ is the integral of its derivative, so if at every value in between every non-AI mode in $\mathcal M_j$ uses at least as many payroll dollars per unit as every AI mode (mode A at its optimal share), then (E1), applied with $\theta=-\tau_L$, shows that $\mathcal B_j$ is nondecreasing in $\tau_L$. The reverse ordering reverses the conclusion, and the same holds for $\vartheta$ with hours in place of payroll dollars.

(c) $\tau_L$ and $\vartheta$ appear in \eqref{eq:ext-effcost}, including mode A's choice of review share, only through $\hat w_j$ and $\hat w^L_j$, and a payroll subsidy of $u_s$ per hour lowers both by $u_s$, the same as $\vartheta=(1-\tau)u_s$. A deductible tax $T_D$ per own hour displaced adds $(1-\tau)T_D(t_j-\bar t_x)$ to the after-tax cost of mode $x$, that is, $T_Dt_j-T_D\bar t_x$ to $\tilde c_x$. The first term is the same for every mode and cancels in $\mathcal B_j$. The second lowers the cost of every own hour by $T_D$, which is the altruism term $-\vartheta\bar t_x/(1-\tau)$ with $\vartheta=(1-\tau)T_D$; in mode A it therefore also moves the optimal review share exactly as that altruism does. \qed

\paragraph{Proposition~\ref{prop:ext-margin}.} (a) $\Sigma^{\rm eco}(\theta)=\sum_jE_j\mathbf 1\{\mathcal B_j(\theta)>0\}/\sum_iE_i$ changes only when some $\mathcal B_j$ changes sign, and then by $E_j/\sum_iE_i$. Each term $v_j(\mathcal B_j(\theta))_+$ is continuous, with derivative $v_jb_j$ where $\mathcal B_j>0$ and zero where $\mathcal B_j<0$, and piecewise affine where $\mathcal B_j$ is affine. Summing gives continuity and \eqref{eq:ext-envelope}.

(b) Let $P_E$ be the spend-weighted distribution of $(\mathcal B_j(\theta_0),b_j)$, normalized to total mass one, and $f(\cdot\mid b)$ the conditional density of $\mathcal B_j(\theta_0)$. Then $\Sigma^{\rm eco}(\theta)=\int\bigl[1-F(-(\theta-\theta_0)b\mid b)\bigr]dP_E(b)$. Assume $|b|$ is integrable and the conditional densities are bounded near zero by a constant. Differentiating under the integral (dominated convergence and the mean value theorem) gives $d\Sigma^{\rm eco}/d\theta|_{\theta_0}=\int b\,f(0\mid b)\,dP_E(b)$. By the definition of conditional expectation for a variable with density, $\int bf(0\mid b)\,dP_E(b)=f_{\mathcal B}(0)\,\mathbb E^E[b\mid\mathcal B=0]$, where $f_{\mathcal B}(0)=\int f(0\mid b)\,dP_E(b)$.

(c) If $\theta$ enters only AI-mode costs affinely, $\tilde c_x(\theta)=\alpha_x+\beta_x\theta$ for $x\in\mathcal M^{\rm AI}$ and the non-AI costs are constant. The minimum of affine functions is concave, so $\mathcal B_j(\theta)$, a constant minus a concave function, is convex. So is $(\mathcal B_j)_+=\max\{\mathcal B_j,0\}$, and a nonnegative combination of convex functions is convex. A convex function lies above its tangent, which gives the last statement. \qed

\paragraph{Proposition~\ref{prop:ext-customer}.} (a) By Remark~\ref{rem:taxes}, the after-tax problem is Part I's problem after the substitutions $m_0\mapsto(1-\tau)m_0$, $k\mapsto(1-\psi\tau)k$ and $\lambda\mapsto(1-\tau)^{-\alpha}\lambda$. A proportional change in $\lambda$ is the same proportional change in the substituted $\lambda$. By Proposition~\ref{prop:envelope}, applied with firm-specific zero-spend growth $G_i$ as in Remark~\ref{rem:assumptions}(b),
\[
dR_X^*=(1-\tau)\,dm_0+\bigl[X-(1-\psi\tau)k\bigr]dG_i+\bigl[X-(1-\psi\tau)k\bigr]\Delta^*\frac{d\lambda}{\lambda}
\]
to first order, and $dG_i=-\delta^G$, $d\lambda/\lambda=-\delta^\lambda$. Let $b^{\rm pec}$ be the pecuniary saving per unit, so that $dm_0=b^{\rm pec}v_j/S$. Then $dR_X^*>0$ if and only if $b^{\rm pec}>[X-(1-\psi\tau)k](\delta^G+\Delta^*\delta^\lambda)S/[(1-\tau)v_j]=\mathcal C_{x,j}$. By Proposition~\ref{prop:dynamic}(c), the maximand of the Hamilton--Jacobi--Bellman equation at share $s$ is the Rule-of-$X$ score with $X=v'(s)$, so the first-order change in that maximand is the change in the firm's own criterion.

(b) With $G_i$ in place of $G$, the small-share multiple is the smallest positive root of $\mathcal Y_0(v_0;m_0,G_i,\lambda):=(r-G_i)v_0-m_0+kG_i-C_\alpha[\lambda(v_0-k)]^{1/(1-\alpha)}=0$; for $G_i=G$ this is \eqref{eq:gordon}. Now $\partial_{v_0}C_\alpha[\lambda(v_0-k)]^{1/(1-\alpha)}=\alpha^{\alpha/(1-\alpha)}\lambda^{1/(1-\alpha)}(v_0-k)^{\alpha/(1-\alpha)}=\Delta^\circ$ by \eqref{eq:Dstar} with $X=v_0$. Hence $\partial\mathcal Y_0/\partial v_0=r-G_i-\Delta^\circ=r-g^\circ$, which is positive by Proposition~\ref{prop:dynamic}(d). Further, $\partial\mathcal Y_0/\partial m_0=-1$ and $\partial\mathcal Y_0/\partial G_i=-(v_0-k)$. Since $\partial_\lambda C_\alpha[\lambda y]^{1/(1-\alpha)}=\alpha^{\alpha/(1-\alpha)}\lambda^{\alpha/(1-\alpha)}y^{1/(1-\alpha)}=y\Delta^\circ/\lambda$, $\partial\mathcal Y_0/\partial\lambda=-(v_0-k)\Delta^\circ/\lambda$. By the implicit function theorem,
\[
dv_0=\frac{dm_0+(v_0-k)\,dG_i+(v_0-k)\Delta^\circ\,d\lambda/\lambda}{r-g^\circ},
\]
which is \eqref{eq:ext-dv0} with $dG_i=-\delta^G$ and $d\lambda/\lambda=-\delta^\lambda$. Its numerator is the score change of (a) with $X=v_0$ and $\tau=0$. \qed

\paragraph{Corollary~\ref{cor:ext-floor}.} By Proposition~\ref{prop:ext-agi}(b), with unbounded capability the review share of mode A for $j\in\mathcal J^{\rm per}$ is $\xi^{\min}(c_j)$ and its errors vanish. Its human cost per unit is therefore $w_j\xi^{\min}(c_j)t^R_j+\mathrm{lic}_jw^L_jt^L_j$, while machine spend belongs to $c_c$. Impermissible activities keep the cost $c_{x^0_j,j}$. Summing over activities and dividing by revenue gives \eqref{eq:ext-amin}. Its part in cost of revenue is a lower bound on $c_h$, and its operating part a lower bound on $a$. Corollary~\ref{cor:ceiling} bounds $R_\infty$ by its value at the smallest admissible $c_h$ and $a$, which gives the stated version. Raising $\xi^{\min}(c)$ by $d\xi$ raises the residual cost by $\sum_{j:c_j=c}v_jw_jt^R_j\,d\xi/S$ over activities whose floor binds, and $R_{\max}$ falls by $(1-\tau)$ times that amount. \qed

\paragraph{Proposition~\ref{prop:ext-shapley}.} (a) Order the groups by a permutation $o$ and let $P_o(k)$ be the groups preceding $k$. Then $\sum_k[f(P_o(k)\cup\{k\})-f(P_o(k))]=f(\{1..n_g\})-f(\emptyset)$, by telescoping. For a uniformly random permutation, $\Pr(P_o(k)=T)=|T|!(n_g-|T|-1)!/n_g!$, so $\mathrm{Sh}_k$ is the expected marginal contribution of $k$. Averaging the telescoping identity over permutations gives (a). (b) If $f(T)=f(\emptyset)+\sum_{i\in T}g_i$, every marginal contribution of $k$ equals $g_k$, and the weights sum to one. (c) The verification script exhibits a parameter scenario in which the one-at-a-time changes sum to 0.069 while the total change is 0.015. \qed

\paragraph{Assumption~\ref{ass:ext-capability}(ii) and (iii).} From $p(d,t)=\bar q$, $(d/H)^{1/\eta_r}=(1-\bar q)/\bar q$, so $d=H[(1-\bar q)/\bar q]^{\eta_r}=H/\chi(\bar q)$. At 80\%, $\chi=4^{\eta_r}$, so a ratio of 4 to 6 gives $\eta_r=\ln\chi(0.8)/\ln4\in[1,\ln6/\ln4]$. For (iii), $d/[d^g+(1-\iota C)d^c]=1/[1-s^{\rm ctx}+(1-\iota C)s^{\rm ctx}]=1/(1-s^{\rm ctx}\iota C)$.

\paragraph{Proposition~\ref{prop:ext-feasible}.} (a) The condition $\mathcal H\ge\bar d_j$, that is $H_0\,2^{(t-L_W\mathrm{ph}_j-t_0)/T_d}\kappa_{\rm ctx}\ge d\chi(\bar q)$, is equivalent to
\[
t\;\ge\;t_0+L_W\mathrm{ph}_j+T_d\log_2\frac{d\chi(\bar q)}{H_0\kappa_{\rm ctx}}.
\] (b) The delay is $T_d[\log_2\chi(\bar q')-\log_2\chi(\bar q)]=\eta_rT_d\log_2\{[\bar q'/(1-\bar q')]/[\bar q/(1-\bar q)]\}$. The numbers follow from $\eta_r=\ln5/\ln4$ and the odds 1, 4, 19 and 99. (c) $n^{\rm dbl}_j$ falls by $\log_2(\kappa'/\kappa)$, and $\kappa_{\rm ctx}\le1/(1-s^{\rm ctx}_j)$ at $\iota_jC_j=1$. (d) By (a), $j$ is feasible at $t$ if and only if $n^{\rm dbl}_j\le(t-t_0-L_W\mathrm{ph}_g)/T_d$. The spend share is the sum over groups of the group weight times $F_g$ evaluated there. Under normality, $F_g(u)=\Phi((u-\bar n^{\rm dbl})/\varsigma_{\rm dbl})$, which is one half at $t=t_0+T_d\bar n^{\rm dbl}+L_W\mathrm{ph}_g$; the 10\% and 90\% dates are $z_{0.9}\varsigma_{\rm dbl}T_d$ either side of it. (e) Replacing $T_d$ by $T_d'$ multiplies $(t-t_0)/T_d$ by $T_d/T_d'$, and multiplying every $\kappa$ by $\kappa'/\kappa$ lowers every $n^{\rm dbl}_j$ by $\log_2(\kappa'/\kappa)$. (f) Under the plateau, $\mathcal H=H(\min\{t-L_W\mathrm{ph}_j,t_p\})\le H(t_p)$, with equality from $t_p+L_W\mathrm{ph}_j$ on, so the test of (a) can hold at some date if and only if $\bar d_j\le H(t_p)$, and when it does it first holds at the date of (a). The share of (d) is then $\sum_gE_gF_g((\min\{t-L_W\mathrm{ph}_g,t_p\}-t_0)/T_d)$, constant from $t_p+L_W$ on. Since $\bar d_j$ is inversely proportional to $\kappa_{{\rm ctx},j}$, multiplying $\kappa_{{\rm ctx},j}$ by $\bar d_j/H(t_p)$ brings $\bar d_j$ down to $H(t_p)$. \qed

\paragraph{Proposition~\ref{prop:ext-energy}.} (a) After the fall, AI spend per unit is $(1-\varepsilon_E)c^{\rm en}+c^{\rm hw}=(1-s_E\varepsilon_E)c^{\rm mach}$. The trend reaches the same level after $\Delta t$ with $e^{-\delta_\pi\Delta t}=1-s_E\varepsilon_E$. (b) With $c^{\rm en}(t)=c^{\rm en}_0e^{-\delta_{\rm en}t}$ and $c^{\rm hw}(t)=c^{\rm hw}_0e^{-\delta_{\rm hw}t}$, $s_E=c^{\rm en}/(c^{\rm en}+c^{\rm hw})$, and differentiating gives $\dot s_E=s_E(1-s_E)(\delta_{\rm hw}-\delta_{\rm en})$. This is a logistic equation, whose solution tends to one if $\delta_{\rm hw}>\delta_{\rm en}$ and to zero if $\delta_{\rm hw}<\delta_{\rm en}$. (c) $\mathcal J^{\rm fea}$ depends on capability and readiness, $\mathcal J^{\rm per}$ on $\Theta$ and feasibility, and $\mathcal J^{\rm pre}$ on costs with the machine terms of AI modes set to zero (Lemma~\ref{lem:ext-nested}(b)). Energy enters none of them. (d) $c_c$ is AI spend per dollar of revenue, so it falls by $s_E\varepsilon_Ec_c$, and $R_{\max}$ of Corollary~\ref{cor:ceiling} rises by $(1-\tau)$ times that amount. Under free entry, Proposition~\ref{prop:entry}(b) gives a long-run score that depends on unit cost only through $b/\hat c$, so the gain in margin does not persist. (e) The formulas do not depend on the activity type. \qed

\paragraph{Proposition~\ref{prop:ext-agi}.} (a) For fixed $d_j$, $(d_j/(\kappa_{{\rm ctx},j}\mathcal H_j))^{1/\eta_r}\to0$ as $H\to\infty$, so $p\to1$, $e_j\to0$ and the capability test holds. Feasibility then reduces to readiness and, for physical activities, the world-model test. (b) By Proposition~\ref{prop:tasks-review}(d), as $e_j\to0$ the review share tends to $\xi^{\min}(c_j)$ and the quality cost to $\xi^{\min}(c_j)\hat w_jt^R_j$, so the error term vanishes. With AI spend per unit tending to its floor $c^{\rm mach,\infty}$, \eqref{eq:ext-effcost} tends to \eqref{eq:ext-cagi}, in which $H$ does not appear. (c) The limit universe is $\{j:\ \mathcal M^{\rm AI}\cap\mathcal M_j^\infty\neq\emptyset,\ \min_0\tilde c>\min_{\rm AI}\tilde c^{\,\infty}\}$, which depends only on the objects listed. It carries all of the spend if and only if every activity with positive weight belongs to it, which is the conjunction of (i) to (iv). (d) This follows from Corollary~\ref{cor:ext-floor}, whose floors do not involve $H$. If a failure rate of Section~\ref{sec:tasks-failure} does not vanish as $H\to\infty$, the corresponding loss of Proposition~\ref{prop:tasks-cerr} is a limit cost that does not involve $H$ either, and it enters the ceiling as that proposition states. \qed

\paragraph{Proposition~\ref{prop:ext-timing}.} Assume $\int_0^\infty e^{-rt}|W(t)|\,dt<\infty$.
(a) The derivative of $\int_T^\infty e^{-rt}W(t)\,dt$ is $-e^{-rT}W(T)$, and that of $e^{-rT}F_0e^{-\delta_FT}$ is $-(r+\delta_F)e^{-rT}F_A(T)$.
(b) $\mathcal O'(T)=-e^{-rT}N(T)$. If $N$ is strictly increasing and $N(0)\ge0$, then $\mathcal O'<0$ for $T>0$, so $T^*=0$. If $N(0)<0$ and $N(\hat T)=0$, $\mathcal O$ increases on $[0,\hat T)$ and decreases afterwards, so $T^*=\hat T$ is the unique maximizer. If $N<0$ throughout, $\mathcal O$ increases and tends to zero, so $\mathcal O<0$ everywhere and adoption never pays. Under pay-as-you-go pricing, future values of $c^{\rm var}_A$ enter $\mathcal O'(T)$ only through $W(T)$.
(c) $\mathcal O(T^*)-\mathcal O(0)=\int_0^{T^*}\mathcal O'(t)\,dt=\int_0^{T^*}e^{-rt}[-N(t)]\,dt$, and $N<0$ on $[0,T^*)$. The example in Section~\ref{sec:ext-dynamics} has $\mathcal O(0)=\bar W/r-F_0=20>0$ and $N(0)=4-8<0$.
(d) (i) $N(T)=\bar W-(r+\delta_F)F_0e^{-\delta_FT}=0$ gives the formula. If it is negative, $N(0)\ge0$ and $T^*=0$. (ii) With $\delta_F=0$, $N(T)=v(c_0-c^Q-\pi_0e^{-\delta_\pi T})-rF_0$, which is increasing. It has a root if and only if $c_0-c^Q>rF_0/v$, and the root is the formula, truncated at zero.
(e) With lock-in, $\mathcal O(T)=\int_T^\infty e^{-rt}v(t)[c_0(t)-c^{\rm var}_A(T)]\,dt-e^{-rT}F_A(T)$. Differentiating,
\[
\mathcal O'(T)=-e^{-rT}W(T)-\dot c^{\rm var}_A(T)\int_T^\infty e^{-rt}v(t)\,dt+(r+\delta_F)e^{-rT}F_A(T),
\]
and $\int_T^\infty e^{-rt}ve^{g_vt}dt=e^{-rT}v(T)/(r-g_v)$.
(f) Write $y(t)=e^{-rt}W(t)$. The first term of \eqref{eq:ext-OT} is $\int y(T+\tau)\varrho(\tau)\,d\tau$ over $\tau\ge0$. Its derivative in $T$ is $\int y'(T+\tau)\varrho(\tau)\,d\tau$. Integrating by parts, with $\varrho(0)=0$ and $\varrho=1$ beyond $T_{\rm par}$, this is
\[
-y(T)\varrho(0)-\int_0^{T_{\rm par}}y(T+\tau)\,d\varrho(\tau)=-e^{-rT}\int_0^{T_{\rm par}}e^{-r\tau}W(T+\tau)\,d\varrho(\tau).
\] The second term contributes $e^{-rT}[r\mathcal T(T)-\mathcal T'(T)]$, which gives \eqref{eq:ext-Opar}. With $W$ constant on the window the integral is $\bar\varrho_rW$. Since $e^{-r\tau}\in[e^{-rT_{\rm par}},1]$ and $\int d\varrho=1$, $\bar\varrho_r\in[e^{-rT_{\rm par}},1]$. For $\varrho(\tau)=\tau/T_{\rm par}$, $\bar\varrho_r=(1-e^{-rT_{\rm par}})/(rT_{\rm par})$. A component $\mathcal T_i=\mathcal T_{i0}e^{-\delta_iT}$ contributes $r\mathcal T_i-\mathcal T_i'=(r+\delta_i)\mathcal T_i$, and a component growing at rate $g_v$ contributes $(r-g_v)\mathcal T_i$. A break fee $b_0(1-T/T_c)_+$ is continuous, and its contribution $rb_0(1-T/T_c)+b_0/T_c$ for $T<T_c$ drops to zero at $T_c$. If the bracket in \eqref{eq:ext-Opar} is positive just before $T_c$ and nonpositive after it, then $\mathcal O$ increases up to $T_c$ and decreases afterwards under the single-crossing condition of (b), so $T^*=T_c$. (g) (i) Let $\Delta W:=W-W_p\ge0$, zero before $t_p$, and $\Delta(T):=\mathcal O(T)-\mathcal O_p(T)=\int_T^\infty e^{-rt}\varrho(t-T)\Delta W(t)\,dt$. For each $t$, $\varrho(t-T)$ is nonincreasing in $T$, and raising the lower limit drops a nonnegative integrand, so $\Delta$ is nonincreasing. For $T\le t_p-T_{\rm par}$ the ramp is complete before $t_p$, so $\Delta(T)=\int_{t_p}^\infty e^{-rt}\Delta W\,dt$; with $T_{\rm par}=0$ this holds for every $T\le t_p$. Hence $\max_T\mathcal O_p(T)-\mathcal O_p(0)=\max_T[\mathcal O(T)-\mathcal O(0)+\Delta(0)-\Delta(T)]\ge\max_T\mathcal O(T)-\mathcal O(0)$. By \eqref{eq:ext-Opar}, $\mathcal O_p'(T)-\mathcal O'(T)=e^{-rT}\int_0^{T_{\rm par}}e^{-r\tau}\Delta W(T+\tau)\,d\varrho(\tau)\ge0$ (in the base case, $e^{-rT}\Delta W(T)$). If $\mathcal O$ increases before its maximizer $T^*$, so does $\mathcal O_p$, and under single crossing its maximizer is at least $T^*$. In the base case, if $N_p(T)=W_p(T)-(r+\delta_F)F_A(T)<0$ at every date, $\mathcal O_p$ increases toward zero and is negative, as in (b). (ii) Under lock-in, $\mathcal O_L(T)=\int_T^\infty e^{-rt}v(t)[c_0(t)-c^{\rm var}_A(T)]\,dt-e^{-rT}\mathcal T(T)$. The plateau leaves $c^{\rm var}_A(T)$ unchanged for $T\le t_p$ and raises it for $T>t_p$, because the quality cost no longer falls while the price part follows the same path. So $\mathcal O_{L,p}=\mathcal O_L$ on $[0,t_p]$ and $\mathcal O_{L,p}\le\mathcal O_L$ after, which gives $\max_T\mathcal O_{L,p}\le\max_T\mathcal O_L$ with $\mathcal O_{L,p}(0)=\mathcal O_L(0)$. After $t_p$, $\dot c^{\rm var}_A$ is the price part alone, which enters the hurdle of (e). (iii) follows from Proposition~\ref{prop:ext-feasible}(f). \qed

\paragraph{Lemma~\ref{lem:ext-retained}.} At a date $t$, firm $i$ maximizes $(P-\hat c_i)x_i$ with $P=A-B\sum_lx_le^{-Gt}$. Write $\tilde x_l=x_le^{-Gt}$. The first-order condition is $P-\hat c_i=B\tilde x_i$. Summing over the $n_c$ firms gives $n_cP-\sum\hat c_l=B\sum\tilde x_l=A-P$, so $P=(A+\sum\hat c_l)/(n_c+1)$. Flow profit is $B\tilde x_i^2e^{Gt}=(P-\hat c_i)^2e^{Gt}/B$. (a) Lowering $n_r+1$ costs by $\delta_c$ lowers $\sum\hat c_l$ by $(n_r+1)\delta_c$ and $P$ by $(n_r+1)\delta_c/(n_c+1)$. A cutter's margin changes by $\delta_c-(n_r+1)\delta_c/(n_c+1)=(n_c-n_r)\delta_c/(n_c+1)$, and a non-cutter's by $-(n_r+1)\delta_c/(n_c+1)$. (b) $\bar\varphi=(n_r+1)/(n_c+1)$, and $[n_c-\mathrm{rep}(n_c-1)]/(n_c+1)\to1-\mathrm{rep}$. Lemma~\ref{lem:prices} defines $\varphi=(dP/du)(u/P)=\bar\varphi u/P$. \qed

\paragraph{Proposition~\ref{prop:ext-should}.} (a) Let $\hat\delta:=\delta_c/(n_c+1)$. By Lemma~\ref{lem:ext-retained}, when $n_r$ rivals adopt, a non-adopter's margin is $\mathrm{mg}-n_r\hat\delta$, and it rises to $\mathrm{mg}+(n_c-n_r)\hat\delta$ if the firm adopts as well. The no-exit condition $\mathrm{mg}\ge(n_c-1)\hat\delta$ keeps every output positive. Then $\mathcal E^{\rm priv}(n_r)=[(\mathrm{mg}+(n_c-n_r)\hat\delta)^2-(\mathrm{mg}-n_r\hat\delta)^2]/B=n_c\hat\delta[2\mathrm{mg}+(n_c-2n_r)\hat\delta]/B$, which is strictly decreasing in $n_r$. Also $\mathcal E^{\rm eq}=\hat\delta(2\mathrm{mg}+\hat\delta)/B$, and
\[
\mathcal E^{\rm priv}(n_c-1)-\mathcal E^{\rm eq}=\frac{\hat\delta}{B}\bigl[n_c(2\mathrm{mg}-(n_c-2)\hat\delta)-(2\mathrm{mg}+\hat\delta)\bigr]=\frac{\hat\delta(n_c-1)}{B}\bigl[2\mathrm{mg}-(n_c-1)\hat\delta\bigr]\ge0.
\]
If $\mathcal K<\mathcal E^{\rm priv}(n_c-1)$, adopting beats not adopting whatever the rivals do, so all adopt. Each firm then gains $\mathcal E^{\rm eq}-\mathcal K$ over the status quo: positive in the value-creating case, nonpositive in the defensive case. If $\mathcal E^{\rm priv}(n_c-1)\le\mathcal K<\mathcal E^{\rm priv}(0)$, let $\hat n_r=\min\{n_r:\mathcal E^{\rm priv}(n_r)\le\mathcal K\}\in\{1,\dots,n_c-1\}$. In a profile with $\hat n_r$ adopters, an adopter, facing $\hat n_r-1$ adopting rivals, gains $\mathcal E^{\rm priv}(\hat n_r-1)>\mathcal K$, and a non-adopter, facing $\hat n_r$, would gain $\mathcal E^{\rm priv}(\hat n_r)\le\mathcal K$. The profile is therefore an equilibrium with partial adoption. If $\mathcal K\ge\mathcal E^{\rm priv}(0)$, not adopting is dominant. Finally, $\mathcal E^{\rm eq}=\hat\delta(2\mathrm{mg}+\hat\delta)/B\to0$ as $\hat\delta\to0$.
(b) By Proposition~\ref{prop:entry}(b), free entry drives the entrants' price to their full unit cost $\hat c_E$. Base customers stay if and only if $P_I\le\hat c_E+p_{\rm sw}$. If $P_I$ exceeds this, keeping the base requires a price of at most $\hat c_E+p_{\rm sw}$, and in this reduced form, with the base's demand fixed below the switching threshold, the incumbent sets exactly that price. The price no longer depends on the incumbent's cost, so $\bar\varphi=0$, and its unit margin is $\hat c_E+p_{\rm sw}-\hat c_I$.
(c) A nonnegative margin at the best attainable cost is necessary and sufficient for survival, and the decomposition holds by definition.
(d) Automating $j$ lowers unit cost by $\mathcal B_jv_j/Y_I$ and changes value by $\mathrm{NPV}_j\le0$. An activity with $\mathcal B_j\le0$ has $\mathrm{NPV}_j\le-\mathcal T_j\le0$ and covers nothing, so dropping it from a feasible cover keeps the cover feasible and does not raise its cost. In the linear relaxation $\min\sum_jc_jy_j$ subject to $\sum_ja_jy_j\ge\mathcal D$ and $0\le y_j\le1$, with $c_j=-\mathrm{NPV}_j\ge0$ and $a_j>0$, order the activities by $c_j/a_j$. Set $y_j=1$ until the next activity would overshoot, and take that one fractionally. With $u$ the ratio $c_j/a_j$ of the fractional activity, the dual solution $(u,z_j=(a_ju-c_j)_+)$ is feasible and satisfies complementary slackness, so the greedy solution is optimal, with at most one fractional variable.
(e) If a defensive activity is not operating by $T_M$, the incumbent's unit cost exceeds $\hat c_E+p_{\rm sw}$ from that date, and by (c) it loses its base. (f) Part (i) is Proposition~\ref{prop:ext-feasible}(f). For (ii), the entrant automates every feasible, permissible and economical activity at $\iota=1$, and its unit cost is nonincreasing in capability, through feasibility and through quality costs. Both are frozen from $t_p$ on, so $\hat c_{E,p}(t)\ge\hat c_E(t)$, with equality before $t_p$. The condition of (b), $P_I>\hat c_E(t)+p_{\rm sw}$, therefore holds on a subset of the dates at which it held, so its first date weakly rises, and if it never holds, (b) to (e) never apply and the defensive set is empty. For (iii), the test of (c), $\hat c_I^*\le\hat c_E+p_{\rm sw}$, has capability on both sides. For (iv), $\mathrm{NPV}_j$ is increasing in the path of $\mathcal B_j$, which the plateau lowers after $t_p$. \qed

\paragraph{Proposition~\ref{prop:ext-team}.} Let $g(n)=n-\chi_hn(n-1)/2$ for $n\ge1$ and $g(n)=n$ for $n\le1$. On $[1,\infty)$, $g$ is strictly concave with $g'(n)=1-\chi_h(n-\tfrac12)$, it is maximized at $\hat n=1/\chi_h+\tfrac12$, and $g(\hat n)=(1/\chi_h+\tfrac12)(\tfrac12+\chi_h/4)=(1+\chi_h/2)^2/(2\chi_h)=\mathcal W_{\max}$.
(b) The minimal team is the smallest $n$ with $g(n)\ge \mathcal W$. For $\mathcal W\le1$ this is $n=\mathcal W$. For $1<\mathcal W\le \mathcal W_{\max}$ it is the smaller root of $\tfrac{\chi_h}2n^2-(1+\tfrac{\chi_h}2)n+\mathcal W=0$, which is \eqref{eq:ext-nstar}; there is no root if $\mathcal W>\mathcal W_{\max}$.
(e) On $(1,\mathcal W_{\max})$, $n^{*\prime}(\mathcal W)=1/g'(n^*)=1/[1-\chi_h(n^*-\tfrac12)]>1$, and $n^{*\prime\prime}=-g''(n^*)\,n^{*\prime}/g'(n^*)^2=\chi_h(n^{*\prime})^3>0$. At $\mathcal W=1$ the slope rises from 1 to $1/(1-\chi_h/2)$, so $n^*$ is convex on $[0,\mathcal W_{\max}]$ with $n^*(0)=0$. Therefore $n^*(\mathcal W)/\mathcal W$ is nondecreasing, and $n^*(\mathcal W/L)\le n^*(\mathcal W)/L$ for $L>1$.
(a) With $\chi_a=0$, $n_a=\nu_a\mathcal W_0/Y_a$, and the cost is $C(\nu_a)=w\,n^*(\mathcal W(\nu_a))+p_a\nu_a\mathcal W_0/Y_a$, where $\mathcal W'(\nu_a)=-\mathcal W_0(1/A_g-\xi^t)$. If $1/A_g\le\xi^t$, agents do not save human time, and $\nu_a=0$ when $p_a>0$. Otherwise $C'(\nu_a)=\mathcal W_0[p_a/Y_a-w(1/A_g-\xi^t)\,n^{*\prime}(\mathcal W(\nu_a))]$. Since $n^{*\prime}\ge1$, the condition $p_a/Y_a<w(1/A_g-\xi^t)$ gives $C'<0$ throughout and $\nu_a=\bar\nu_a$. Since $n^*$ is convex and $\mathcal W$ is decreasing in $\nu_a$, $C'$ is nondecreasing, so $C$ is convex. Its minimizer is zero if $C'(0)\ge0$, the boundary $\bar\nu_a$ if $C'(\bar\nu_a)\le0$, and otherwise the interior point where $w\,n^{*\prime}(1/A_g-\xi^t)=p_a/Y_a$, with $w\,n^{*\prime}=\tilde w$.
(c) $g$ is increasing on $[0,\hat n]$ and $g(1)=1$, so $n^*\le1$ if and only if $\mathcal W\le1$, that is, $\mathcal W_0\le\Lambda$.
(d) Both terms of $\Lambda^{-1}$ are nonnegative, so $\Lambda^{-1}\ge(1-\nu_a)/A_g$ and $\Lambda^{-1}\ge\nu_a\xi^t$. As $\nu_a\to1$, $\Lambda^{-1}\to\xi^t$, and $\xi^t\ge\xi^{\min}t^R/t$ gives the floor. $\Lambda\ge10$ if and only if $(1-\nu_a)/A_g+\nu_a\xi^t\le0.1$. With $A_g\le2$ this needs $(1-\nu_a)/2\le0.1$, so $\nu_a\ge0.8$, and then $\xi^t\le0.1/\nu_a\le0.125$, with strict inequality unless $\nu_a=1$.
(f) Profit $P[nA_g-\chi_hn(n-1)/2]-wn$ is strictly concave in $n$, and its first-order condition $PA_g-P\chi_h(n-\tfrac12)-w=0$ gives $n^\pi$, increasing in $A_g$.
(g) The algebra of (b) applies with $(Y_a,\chi_a)$ in place of $(1,\chi_h)$. \qed

\paragraph{Proposition~\ref{prop:ext-hierarchy}.} Each worker escalates the share $e^{-z_w}$ of problems, each of which takes $\kappa_c$ of a manager's unit of time, so one manager serves $1/(\kappa_ce^{-z_w})$ workers. The cost per worker is therefore $w_w+c_Kz_w+\kappa_ce^{-z_w}(w_m+c_Kz_m)$.
(a) The cost is strictly convex in $z_w$, and the first-order condition $c_K=\kappa_ce^{-z_w}(w_m+c_Kz_m)$ gives $z_w^*$ and the span.
(b) $\partial\mathrm{sp}^*/\partial c_K=-w_m/c_K^2<0$, and $\partial z^*_w/\partial c_K=z_m/(w_m+c_Kz_m)-1/c_K=-w_m/[c_K(w_m+c_Kz_m)]<0$. Also $\partial z^*_w/\partial\ln\kappa_c=1$, and $\mathrm{sp}^*$ does not contain $\kappa_c$. For a general density $f$ with survivor function $\bar F$ and $\mathcal K_m=w_m+c_Kz_m$, the first-order condition is $\kappa_c\mathcal K_mf(z^*)=c_K$, and the second-order condition is $f'(z^*)<0$. Differentiating in $\ln\kappa_c$ gives $dz^*/d\ln\kappa_c=-f/f'$. Since $\mathrm{sp}^*=1/(\kappa_c\bar F(z^*))$, $d\ln\mathrm{sp}^*/d\ln\kappa_c=-1+(f/\bar F)(-f/f')$. With hazard $\lambda_z=f/\bar F$, $\ln f=\ln\lambda_z+\ln\bar F$, so $-f'/f=\lambda_z-\lambda_z'/\lambda_z>0$. Hence $d\ln\mathrm{sp}^*/d\ln\kappa_c=-1+\lambda_z/(\lambda_z-\lambda_z'/\lambda_z)$, which is negative if $\lambda_z'<0$, zero if $\lambda_z'=0$ and positive if $\lambda_z'>0$.
(c) By memorylessness, $z-z_A$ given $z>z_A$ is distributed ${\rm Exp}(1)$. Knowledge below $z_A$ is never used by humans, so a human worker who knows $[z_A,z_A+\zeta]$ pays $c_K\zeta$, and a manager pays $c_K(z_m-z_A)$. The cost per human worker is the expression of (a) with $z_m$ replaced by $z_m-z_A$, so the optimum and the span follow from (a). Only the share $e^{-z_A}$ of problems reaches humans, so production workers and managers per problem both fall by that factor.
(d) If $d$ is Pareto with $\Pr(d>x)=(x/d_{\min})^{-\zeta_P}$, then $\Pr(\zeta_P\ln(d/d_{\min})>u)=e^{-u}$. The AI layer solves $d\le H$, that is, $z\le z_A=\zeta_P\ln(H/d_{\min})$, so the human share is $(H/d_{\min})^{-\zeta_P}$, and with $H$ doubling every $T_d$ it falls at the rate $\zeta_P\ln2/T_d$. With span $\mathrm{sp}^*$ at every level, the number of levels above $N_w$ workers is about $\ln N_w/\ln\mathrm{sp}^*$, and multiplying $N_w$ by $e^{-z_A}$ removes $z_A/\ln\mathrm{sp}^*$ of them. \qed

\paragraph{Lemma~\ref{lem:ext-staffing}.} (a) For each activity, $v_jt_j=v_j\bar t_{x^*_j}+v_j(t_j-\bar t_{x^*_j})$. Divide by $t_{\rm FTE}$ and sum; in mode O, part of the displaced hours moves to the vendor. (b) During the parallel run of Section~\ref{sec:tasks}, the legacy mode carries the share $1-\varrho(\tau)$ of volume, the new mode carries $\varrho(\tau)$, legacy capacity is held at the extra share $\omega_{\rm par}$, and reconciliation takes $\nu_{\rm rec}$ of the in-house hours. In-house FTE for process $p$ is therefore $(1-\varrho)\mathrm{FTE}^0_p+\varrho\,\mathrm{FTE}^1_p+\omega_{\rm par}\mathrm{FTE}^0_p+\nu_{\rm rec}\mathrm{FTE}^{\mathrm H}_p$. Summing over processes and adding the transformation team gives the expression, and setting $\tau=0$ gives the hump. (c) The pace cannot exceed what the organization absorbs, $\kappa_{\rm abs}\mathrm{FTE}_{\rm base}$, nor what its change managers can support, $N_{\rm CM}\mathrm{FTE}^{\rm CM}$. \qed

\paragraph{Proposition~\ref{prop:ext-waves}.} (a) Let $p$ be in wave $k$ and $q$ in a later wave $k'$, with $W_p<W_q$. Swapping them leaves every wave's count, and hence $C^{\rm dr}$ and $O_w$, unchanged, and changes the cost by $T_{\rm par}[(kW_q+k'W_p)-(kW_p+k'W_q)]=T_{\rm par}(k'-k)(W_p-W_q)<0$. An optimal plan therefore has nonincreasing savings across waves. In the fractional relaxation with capacity weights and no congestion, let $p$ have mass in a later wave $k'$ than some mass of $q$ in wave $k$, with $W_p/a_p>W_q/a_q$. Moving $\epsilon/a_p$ of $p$ to wave $k$ and $\epsilon/a_q$ of $q$ to wave $k'$ keeps both capacities, and changes the cost by $T_{\rm par}(k'-k)\epsilon(W_q/a_q-W_p/a_p)<0$.
(b) For fixed $n_w$, the problem $\min\sum_k[T_{\rm par}(Wkn_k+c^{\rm dr}_1n_k+c^{\rm dr}_2n_k^2)]$ subject to $\sum n_k=N$ is strictly convex. Its first-order condition $T_{\rm par}(Wk+c^{\rm dr}_1+2c^{\rm dr}_2n_k)=\lambda$ makes $n_k$ affine in $k$ with slope $-W/(2c^{\rm dr}_2)$, and $\sum_kn_k=N$ gives \eqref{eq:ext-front}. With $V_{n_w}:=\sum_k(k-\tfrac{n_w+1}2)^2=n_w(n_w^2-1)/12$, $\sum_kkn_k=N(n_w+1)/2-WV_{n_w}/(2c^{\rm dr}_2)$ and $\sum_kn_k^2=N^2/n_w+W^2V_{n_w}/(4(c^{\rm dr}_2)^2)$. Substituting, $T_{\rm par}[WN(n_w+1)/2-W^2V_{n_w}/(2c^{\rm dr}_2)+c^{\rm dr}_1N+c^{\rm dr}_2N^2/n_w+W^2V_{n_w}/(4c^{\rm dr}_2)]+n_wO_w$, which is the stated $\mathcal C(n_w)$.
(c) With $n_k=n=N/n_w$, $\mathcal C=T_{\rm par}[WN^2/(2n)+WN/2+c^{\rm dr}_1N+c^{\rm dr}_2Nn]+NO_w/n$. This is strictly convex in $n>0$, and its first-order condition gives $n^{*2}=(WN/2+O_w/T_{\rm par})/c^{\rm dr}_2$. For large $N$, $n^*\approx\sqrt{WN/(2c^{\rm dr}_2)}$, so $n_w^*=N/n^*$ grows like $\sqrt N$.
(d) Add the constraints $n_k\le n_{\rm cap}$, with multipliers $\zeta_k\ge0$. The Kuhn--Tucker conditions are $T_{\rm par}(Wk+C^{\rm dr\prime}(n_k))+\zeta_k=\lambda$. Where the cap does not bind, $C^{\rm dr\prime}(n_k)=\lambda/T_{\rm par}-Wk$, which falls with $k$. The waves that want the most, the earliest, are the ones capped, and the uncapped waves keep the decrement of (b). Every plan has $n_w\ge N/n_{\rm cap}$ waves, so it lasts at least $T_{\rm par}N/n_{\rm cap}=Na/K_{\rm chg}$. Where some $\zeta_k>0$, the envelope theorem gives $\partial\mathcal C^*/\partial n_{\rm cap}=-\sum_k\zeta_k<0$.
(e) With $c^{\rm dr}_2=0$, any plan costs $T_{\rm par}W\sum_kkn_k+T_{\rm par}c^{\rm dr}_1N+n_wO_w\ge T_{\rm par}WN+T_{\rm par}c^{\rm dr}_1N+O_w$, the cost of a single wave, because $\sum_kkn_k\ge N$ and $n_w\ge1$.
(f) During wave $k$, completed processes save $\sum_{i<k}n_iW$ and the wave costs $C^{\rm dr}(n_k)$. With nonincreasing sizes the first term rises and the second falls, so the margin change rises with $k$, and the dip is largest in wave 1. A larger $n_1$ raises $C^{\rm dr}(n_1)$ and lowers $n_w$. \qed

\paragraph{Proposition~\ref{prop:ext-portfolio}.} (a) Taking process 1 first gives $\varpi_1V_1+\delta_dV_2[\varpi_1(\varpi_2+\Delta\varpi_2)+(1-\varpi_1)\varpi_2]=\varpi_1V_1+\delta_dV_2(\varpi_2+\varpi_1\Delta\varpi_2)$. The reverse order is symmetric, and subtracting gives \eqref{eq:ext-order}. A success raises experience to $\eta_{\rm tr}$, so the logit of the other process's factor $\varpi^c_j$ rises by $u:=\upsilon_L\ln(1+\eta_{\rm tr})$. Write $\varpi_j=R_j\varpi^c_j$, where $R_j$ collects the activity and congestion factors of \eqref{eq:tasks-success}, which learning leaves unchanged. Since $d\varpi^c/d\,{\rm logit}\,\varpi^c=\varpi^c(1-\varpi^c)$, $\Delta\varpi_j=R_j\Delta\varpi^c_j=\varpi_j(1-\varpi^c_j)u$ to first order. If $R_1=R_2=R$, then $\varpi_1\Delta\varpi_2-\varpi_2\Delta\varpi_1=R^2\varpi^c_1\varpi^c_2[\rho_u(z_2)-\rho_u(z_1)]$ exactly, where $z_j={\rm logit}\,\varpi^c_j$ and $\rho_u(z):=\operatorname{expit}(z+u)/\operatorname{expit}(z)-1$ is the relative gain. Because $\ln\operatorname{expit}$ is strictly concave, $\ln\operatorname{expit}(z+u)-\ln\operatorname{expit}(z)$ is strictly decreasing in $z$ for $u>0$, so $\rho_u$ is strictly decreasing and the learning term has the sign of $z_1-z_2$, which is the sign of $\varpi_1-\varpi_2$. With $V_1=V_2$, the first term of \eqref{eq:ext-order} has the same sign. Under the recursion, a success in wave $k_0$ leaves $E=\eta_{\rm tr}(1-\delta_E)^{k-1}$ in wave $k_0+k$.
(b) After a first-wave success the platform is known to be sound. Wave 2 is then worth $EV_2(1)=\varpi V-(1-\varpi)C$, assumed positive. After a failure the posterior is $p'$, by Bayes' rule, and the firm continues only if $EV_2(p')>0$. Hence $EV^{\rm seq}=EV_1+\delta_d[p_P\varpi\,EV_2(1)+(1-p_P\varpi)(EV_2(p'))_+]$. By the law of total probability, $p_P\varpi\,EV_2(1)+(1-p_P\varpi)EV_2(p')=EV_1$: the coefficient of $V$ is $\varpi[p_P\varpi+p_P(1-\varpi)]=p_P\varpi$, and that of $C$ is $-(1-p_P\varpi)$. So $EV^{\rm seq}=EV_1+\delta_dEV_1+\delta_d(1-p_P\varpi)[(EV_2(p'))_+-EV_2(p')]$, and $EV^{\rm sim}=2EV_1$, which gives \eqref{eq:ext-seq}.
(c) Let $N^{\rm s}$ be the number of successes. By the law of total variance, $\mathrm{Var}(N^{\rm s})=\mathbb E[\mathrm{Var}(N^{\rm s}\mid\text{platform})]+\mathrm{Var}(\mathbb E[N^{\rm s}\mid\text{platform}])=p_Pn_w\varpi(1-\varpi)+(n_w\varpi)^2p_P(1-p_P)$. No wave succeeds if the platform fails.
(d) Let $X_1,X_2\ge0$ be independent with $\Pr(X>x)=x^{-\zeta_T}\mathcal L(x)$, where $\mathcal L$ is slowly varying. Lower bound: $\Pr(X_1+X_2>x)\ge\Pr(\max X_i>x)=2\Pr(X>x)-\Pr(X>x)^2$. Upper bound: for $\epsilon\in(0,\tfrac12)$, $\{X_1+X_2>x\}\subseteq\{X_1>(1-\epsilon)x\}\cup\{X_2>(1-\epsilon)x\}\cup\{X_1>\epsilon x,X_2>\epsilon x\}$. Hence
\[
\frac{\Pr(X_1+X_2>x)}{\Pr(X>x)}\le2(1-\epsilon)^{-\zeta_T}\frac{\mathcal L((1-\epsilon)x)}{\mathcal L(x)}+\epsilon^{-2\zeta_T}x^{-\zeta_T}\frac{\mathcal L(\epsilon x)^2}{\mathcal L(x)}\;\to\;2(1-\epsilon)^{-\zeta_T}.
\]
Letting $\epsilon\to0$ gives the ratio 2. The same argument with tails $c_1x^{-\zeta_T}\mathcal L$ and $c_2x^{-\zeta_T}\mathcal L$ gives $c_1+c_2$, and induction gives $n_w$. $\mathbb EX^2=\int_0^\infty2x\Pr(X>x)\,dx<\infty$ for $\zeta_T>2$, and $\mathbb EX^3=\int3x^2\Pr(X>x)\,dx=\infty$ for $\zeta_T<3$. If each $X_k$ is capped, the total is at most $n_w$ times the cap.
(e) Let $N^{\rm s}_m$ be the successes among provider $m$'s waves. The $N^{\rm s}_m$ are independent across providers, and by (c) applied to each provider, $\mathrm{Var}(N^{\rm s}_m)=p_Pn^{\rm w}_m\varpi(1-\varpi)+(n^{\rm w}_m)^2\varpi^2p_P(1-p_P)$. Summing over $m$ and using $\sum_m(n^{\rm w}_m)^2=n_w^2H_n$ gives the variance. No wave succeeds only if no wave of any provider does, which has probability $\prod_m\Pr(N^{\rm s}_m=0)\ge\prod_m(1-p_P)=(1-p_P)^M$. \qed

\paragraph{Proposition~\ref{prop:ext-commonmode}.} (a) Events at provider $m$ form a Poisson process of rate $\nu^{\rm cm}$, independent across providers, and each costs $c_m:=T_{\rm con}[X_m+X_m^2/(2Y_f)]$. The annual cost $\sum_mc_mN_m$, with $N_m$ independent Poisson variables of mean $\nu^{\rm cm}$, has mean $\nu^{\rm cm}\sum_mc_m=\nu^{\rm cm}T_{\rm con}[X+HX^2/(2Y_f)]$ and variance $\nu^{\rm cm}\sum_mc_m^2$. Both are sums $\sum_mf_i(X_m)$ of a convex function of the provider exposures, $f_1(y)=y+y^2/(2Y_f)$ and $f_2(y)=[y+y^2/(2Y_f)]^2$, so both are Schur-convex: they rise when the vector of exposures becomes more unequal in the sense of majorization, with $X$ fixed, by the Hardy--Littlewood--P\'olya inequality. The vector $(X,0,\dots,0)$ majorizes every other vector with the same total, so a single provider maximizes both.
(b) With $M$ equal shares, $H=1/M$, which gives the saving. The variance is
\[
\nu^{\rm cm}T_{\rm con}^2M\Bigl[\frac XM+\frac{X^2}{2Y_fM^2}\Bigr]^2=\frac{\nu^{\rm cm}T_{\rm con}^2}M\Bigl[X+\frac{X^2}{2Y_fM}\Bigr]^2\le\frac{\nu^{\rm cm}T_{\rm con}^2}M\Bigl[X+\frac{X^2}{2Y_f}\Bigr]^2.
\] With failover an event at $m$ disrupts $\kappa_{\rm fo}X_m$ and costs $T_{\rm con}[\kappa_{\rm fo}X_m+\kappa_{\rm fo}^2X_m^2/(2Y_f)]$; summing over providers and subtracting from (a) gives the stated value.
(c) $J(M+1)-J(M)=F_M+C^{\rm inf}[(M+1)^{b_v}-M^{b_v}]-\kappa_J/[M(M+1)]$ with $\kappa_J:=\nu^{\rm cm}T_{\rm con}X^2/(2Y_f)$, so adding a provider lowers the cost if and only if $\kappa_J>g(M):=M(M+1)[F_M+C^{\rm inf}((M+1)^{b_v}-M^{b_v})]$. The function $g$ is strictly increasing: $M(M+1)F_M$ is, and $M(M+1)[(M+1)^{b_v}-M^{b_v}]=\int_0^1b_vM(M+1)(M+u)^{b_v-1}du$, whose integrand has log-derivative $1/M+1/(M+1)-(1-b_v)/(M+u)>0$ in $M$ because $1/M\ge1/(M+u)$. Hence $J$ falls while $\kappa_J>g(M)$ and does not fall afterwards, and the first $M$ with $\kappa_J\le g(M)$ minimizes $J$, which is \eqref{eq:ext-providers}. At $M=1$ the condition is the stated one. Without discounts $g(M)=M(M+1)F_M$, and the first $M$ with $M(M+1)\ge(M^\dagger)^2$ is the smallest integer at or above $\rho_J:=[\sqrt{1+4(M^\dagger)^2}-1]/2$. Since $M^\dagger-\tfrac12<\rho_J<M^\dagger$, that integer is $\lfloor M^\dagger\rfloor$ or $\lceil M^\dagger\rceil$ (or one if $M^\dagger<1$).
(d) The expected value is a constant plus $p_gB_{\rm up}\delta_c\eta_c-(1-p_g)T_{\rm con}[\eta_cX+\eta_c^2X^2/(2Y_f)]$, a concave quadratic in $\eta_c$. Its derivative $p_gB_{\rm up}\delta_c-(1-p_g)T_{\rm con}X(1+\eta_cX/Y_f)$ vanishes at the bracketed expression, and truncating to $[\eta_{\min},1]$ gives the maximizer of a concave function on an interval. The derivative at $\eta_c=1$ is nonnegative exactly under the stated condition, and then $\eta_c^*=1$. A regression that lasts $T_{\rm reb}$ instead of $T_{\rm con}$ costs $(1-p_g)(T_{\rm reb}-T_{\rm con})[\eta_cX+(\eta_cX)^2/(2Y_f)]$ more in expectation, which is the value of the rollback path.
(e) The mean of (a) is linear in $T_{\rm con}=T_{\rm mon}+T_{\rm resp}$, which gives the value of the kill switch, and its cross-partial in $H$ and $T_{\rm con}$ is $\nu^{\rm cm}X^2/(2Y_f)>0$. A unit with exposure $x_u$ on provider $m$ bears $\nu^{\rm cm}T_{\rm con}x_u[1+X_m/(2Y_f)]$ a year, with marginal cost $\nu^{\rm cm}T_{\rm con}[1+(X_m+x_u)/(2Y_f)]$, while the firm's marginal cost is $\nu^{\rm cm}T_{\rm con}(1+X_m/Y_f)$; the difference is $\nu^{\rm cm}T_{\rm con}(X_m-x_u)/(2Y_f)\ge0$. If each unit's benefit from exposure is concave, a unit equates its marginal benefit to a smaller marginal cost than the firm would, so in a symmetric equilibrium with more than one unit on the provider, exposure exceeds the firm's optimum. Adding the difference, evaluated at the firm's optimum, as a charge per unit of exposure makes each unit's first-order condition at the optimum coincide with the firm's, and the unit's problem is concave, so the optimum is each unit's best response. Moving the decision to the level that bears the whole cost removes the wedge directly. As $x_u/X_m\to0$ the wedge tends to $\nu^{\rm cm}T_{\rm con}X_m/(2Y_f)$.
(f) With $\nu^{\rm cm}=\nu^{\rm out}+\nu^{\rm chg}(\kappa_{\rm itgc})$ and $\nu^{\rm out}$ independent of $\kappa_{\rm itgc}$, the mean of (a) has derivative $\nu^{\rm chg\prime}(\kappa_{\rm itgc})T_{\rm con}[X+HX^2/(2Y_f)]\le0$ in $\kappa_{\rm itgc}$, which is the stated fall. If control quality costs $C_\kappa(\kappa_{\rm itgc})$, increasing and convex, and $\nu^{\rm chg}$ is convex, the total $\nu^{\rm cm}T_{\rm con}[X+HX^2/(2Y_f)]+C_\kappa$ is convex in $\kappa_{\rm itgc}$ and is minimized where $C_\kappa'=-\nu^{\rm chg\prime}T_{\rm con}[X+HX^2/(2Y_f)]$, so quality is worth raising while its marginal cost is below that amount. The marginal value $-\nu^{\rm chg\prime}T_{\rm con}[X+HX^2/(2Y_f)]$ has derivatives $-\nu^{\rm chg\prime}T_{\rm con}X^2/(2Y_f)\ge0$ in $H$ and $-\nu^{\rm chg\prime}[X+HX^2/(2Y_f)]\ge0$ in $T_{\rm con}$, which gives the comparative statics. Staged release and rollback, part (d), lower the expected cost of a regression given that one occurs, and so lower $\nu^{\rm chg}$ measured in full-exposure events. The link to Proposition~\ref{prop:acct-amin} is Assumption~\ref{ass:acct-layers}, in which the same quality $\kappa_{\rm itgc}$ lowers the common-mode probability of AI control layers. \qed

\paragraph{Proposition~\ref{prop:ext-bandwidth}.} (a) With $\varkappa=ne/\Phi$, gross value is $(\Phi/e)\varpi^0V\,\varkappa Q(\varkappa)$, and $d[\varkappa Q]/d\varkappa=Q+\varkappa Q'=Q(1-\varepsilon_Q)$. A single-peaked $\varkappa Q$ therefore peaks where $\varepsilon_Q=1$ and decreases beyond. Net value has the first-order condition $\varpi^0V[Q+\varkappa Q']=C^{\rm dr\prime}(n)\ge0$, so the optimum lies where $\varkappa Q$ is still nondecreasing, at or below $\hat \varkappa$. For $Q=e^{-c/(1-\varkappa)}$, $\varepsilon_Q=c\varkappa/(1-\varkappa)^2$, so $\varepsilon_Q=1$ is $\varkappa^2-(2+c)\varkappa+1=0$, whose root in $(0,1)$ is $\hat \varkappa=1+c/2-\sqrt{c+c^2/4}$. For $q=e^{-\eta_q(\varkappa-1)_+}$, $\varkappa q=\varkappa$ increases on $[0,1]$. For $\varkappa>1$, $\varepsilon_q=\eta_q\varkappa$, which equals one at $1/\eta_q>1$ if $\eta_q<1$, and exceeds one throughout if $\eta_q\ge1$.
(b) Concurrency is at most $\hat \varkappa\Phi/e$ and each initiative is active for $T_{\rm par}$, so throughput is at most $\hat \varkappa\Phi/(eT_{\rm par})$. Combine with Proposition~\ref{prop:ext-waves}(d).
(c) By the envelope theorem, $dEV^*/d\Phi=\partial_\Phi[n\varpi^0VQ(ne/\Phi)]_{n^*}=n^*\varpi^0VQ'(\varkappa^*)(-\varkappa^*/\Phi)=[n^*\varpi^0VQ(\varkappa^*)/\Phi]\varepsilon_Q(\varkappa^*)$.
(d) In initiatives, capacity is $\Phi/e$. Replacing $e$ by $e^{\rm appr}+e^{\rm prep}/(1+A_\Phi)$ multiplies it by the stated factor, which increases in $A_\Phi$ towards $(e^{\rm appr}+e^{\rm prep})/e^{\rm appr}$ without reaching it.
(e) To first order, delegating frees $e_i$ units of top-level bandwidth, worth $p_\Phi e_i$, and lowers expected value by $(\varpi^{\rm top}_i-\varpi^{\rm del}_i)V_i$.
(f) $n\varpi^0VQ(ne/\Phi)-C^{\rm dr}(n)=n\varpi^0V-\tilde C^{\rm dr}(n)$. The function $n[1-Q(ne/\Phi)]$ is a positive multiple of $\varkappa-\varkappa Q(\varkappa)$, which is convex wherever $\varkappa Q$ is concave, and then $\tilde C^{\rm dr}$ is convex. Proposition~\ref{prop:ext-waves} uses only the convexity and monotonicity of $C^{\rm dr}$. \qed

\paragraph{Remark~\ref{rem:ext-seccap}.} (a) With $\nu=n/T_{\rm par}$, $\ln Q_{\rm sec}=\ln(\Phi_{\rm sec}-\nu)-\ln(\Phi_{\rm sec}-\nu+r_d)$, so
\[
-\frac{d\ln Q_{\rm sec}}{d\ln\nu}=\nu\Bigl[\frac1{\Phi_{\rm sec}-\nu}-\frac1{\Phi_{\rm sec}-\nu+r_d}\Bigr]=\frac{\nu r_d}{(\Phi_{\rm sec}-\nu)(\Phi_{\rm sec}-\nu+r_d)},
\]
and $d\ln\nu=d\ln n$. Its numerator increases and its denominator decreases in $\nu$, so $\varepsilon_{\rm sec}$ increases, from zero at $n=0$ to infinity as $\nu\to\Phi_{\rm sec}$. Since $\varkappa=ne/\Phi$ is proportional to $n$, $d\ln Q(ne/\Phi)/d\ln n=-\varepsilon_Q(\varkappa)$. Hence $d\ln[nQ\,Q_{\rm sec}]/d\ln n=1-\varepsilon_Q-\varepsilon_{\rm sec}$. On the domain where $Q$ and $Q_{\rm sec}$ are positive, the gross value $G(n):=nQ\,Q_{\rm sec}$ is positive, and it tends to zero at the right end of the domain, where $Q_{\rm sec}$ (at $\nu=\Phi_{\rm sec}$) or $Q$ vanishes. Its log-derivative is therefore negative somewhere, so the sum of the elasticities exceeds one somewhere. If the sum is continuous, strictly increasing and below one near $n=0$, it crosses one exactly once, at $\hat n_2$, and $G$ increases before $\hat n_2$ and decreases after it; if the sum starts at or above one, $G$ decreases throughout and $\hat n_2=0$. At $n=\hat\varkappa\Phi/e$ the sum is $1+\varepsilon_{\rm sec}>1$, so $\hat n_2<\hat\varkappa\Phi/e$. The security-only breakpoint solves $\nu r_d=(\Phi_{\rm sec}-\nu)(\Phi_{\rm sec}-\nu+r_d)$; with $y=\Phi_{\rm sec}-\nu$ this is $y^2+2r_dy-r_d\Phi_{\rm sec}=0$, whose positive root $y=\sqrt{r_d^2+r_d\Phi_{\rm sec}}-r_d$ gives $\hat\nu_{\rm sec}$, and $y>0$ gives $\hat\nu_{\rm sec}<\Phi_{\rm sec}$. At $n=T_{\rm par}\hat\nu_{\rm sec}$ the sum is $1+\varepsilon_Q\ge1$, so $\hat n_2\le T_{\rm par}\hat\nu_{\rm sec}$, strictly if $\varepsilon_Q>0$ there. If either single-capacity breakpoint lies outside the domain $\{n:ne<\Phi,\,n<T_{\rm par}\Phi_{\rm sec}\}$, the corresponding bound holds because $\hat n_2$ lies inside it. Beyond $\hat n_2$ another concurrent initiative lowers the gross value and raises $C^{\rm dr}$, so concurrency at an optimum is at most $\hat n_2$; each initiative is active for $T_{\rm par}$, so at most $\hat n_2/T_{\rm par}$ processes are migrated a year, which gives the pace bound together with Proposition~\ref{prop:ext-waves}(d). Finally, $\nu Q_{\rm sec}$, the value of the deployment flow per unit of $V_{\rm dep}$, has log-derivative $(1-\varepsilon_{\rm sec})/\nu$ in $\nu$, so it peaks at $\hat\nu_{\rm sec}$.
(b) The problem is $\max_nEV(n;\Phi,\Phi_{\rm sec})$ subject to $n\le K_{\rm chg}T_{\rm par}/a$, with Lagrangian $EV+\zeta(K_{\rm chg}T_{\rm par}/a-n)$. By the envelope theorem the derivative of the optimal value in a parameter is the partial derivative of the Lagrangian at $n^*$. For $\Phi$, $\partial_\Phi[n\varpi^0VQ(ne/\Phi)Q_{\rm sec}]=n\varpi^0VQ_{\rm sec}Q'(\varkappa)(-\varkappa/\Phi)=[n\varpi^0VQ\,Q_{\rm sec}/\Phi]\varepsilon_Q$. For $\Phi_{\rm sec}$, $\partial Q_{\rm sec}/\partial\Phi_{\rm sec}=r_d/(\Phi_{\rm sec}-\nu+r_d)^2$, so $T_{\rm par}^{-1}\partial EV/\partial\Phi_{\rm sec}=(n/T_{\rm par})\varpi^0VQ\,r_d/(\Phi_{\rm sec}-\nu+r_d)^2=\nu V_{\rm dep}r_d/(\Phi_{\rm sec}-\nu+r_d)^2$ with $V_{\rm dep}=\varpi^0VQ(\varkappa)$, which is \eqref{eq:tasks-secvalue}. A cohort of initiatives turns over every $T_{\rm par}$ years, so $EV/T_{\rm par}$ is value per year, and maximizing $EV/T_{\rm par}-p_{\rm sec}\Phi_{\rm sec}$ over $\Phi_{\rm sec}$ at fixed $n^*$ gives $\nu V_{\rm dep}r_d/(\Phi_{\rm sec}-\nu+r_d)^2=p_{\rm sec}$, the first-order condition of Proposition~\ref{prop:tasks-secreview}(c). For $K_{\rm chg}$ the derivative is $\zeta T_{\rm par}/a$. The Kuhn--Tucker condition is $\partial EV/\partial n=\zeta$, and $\partial EV/\partial n=\varpi^0VQ\,Q_{\rm sec}(1-\varepsilon_Q-\varepsilon_{\rm sec})-C^{\rm dr\prime}(n)$ by the computation in (a); $\zeta>0$ only if the cap binds. The first two shadow values are positive whenever $\varepsilon_Q(\varkappa^*)>0$, which holds when latency or quality congestion rises with load, whatever $\zeta$ is.
(c) Fix $n_w$. The relaxed problem minimizes the convex cost $\sum_kT_{\rm par}[Wkn_k+C^{\rm dr}(n_k)]+n_wO_w$ over $\{(n_k):\sum_kn_k=N,\,0\le n_k\le\bar n\}$. The set of pairs $((n_k),\bar n)$ that satisfy these constraints is convex, so the minimum $\mathcal C^*(\bar n)$, a partial minimum of a jointly convex function, is convex in $\bar n$; it is nonincreasing because the feasible set grows with $\bar n$. Where it is differentiable, the envelope theorem applied to the Lagrangian of the proof of Proposition~\ref{prop:ext-waves}(d) gives $\mathcal C^{*\prime}(\bar n)=-\sum_k\zeta_k$. Now $\bar n$ is the minimum of three functions, each increasing and linear in one capacity, with slopes $T_{\rm par}/a$, $1/e$ and $T_{\rm par}$. If one capacity attains the minimum alone, small changes in the other two leave $\bar n$ unchanged, so they have no value, and by the chain rule the binding one is worth $\sum_k\zeta_k$ times its slope. At a tie, raising one capacity alone leaves $\bar n$ unchanged. \qed

\paragraph{Proposition~\ref{prop:ext-shocks}.} Let $\tau_D\sim{\rm Exp}(\nu_D)$ be the first disruption. A disruption after $T_{\rm emb}$ does not matter, so $V=-(F+I)+\mathbb E[V^{\rm pay}(\tau_D)]$, with $V^{\rm pay}(\tau)=V^{\rm pay}(\infty)$ for $\tau\ge T_{\rm emb}$.
(a) Savings accrue at $t\ge T_{\rm par}$ if no disruption has occurred before $\min\{t,T_{\rm emb}\}$, with probability $e^{-\nu_D\min\{t,T_{\rm emb}\}}$. Their value is therefore $W[\int_{T_{\rm par}}^{T_{\rm emb}}e^{-r_Dt}dt+e^{-\nu_DT_{\rm emb}}\int_{T_{\rm emb}}^\infty e^{-rt}dt]$. The recovered value is $\int_0^{T_{\rm emb}}\nu_De^{-\nu_D\tau}e^{-r\tau}V^{\rm rec}d\tau$. Together these give \eqref{eq:ext-Vshock}. On $(0,T_{\rm par})$, $V^{\rm pay}(\tau)=e^{-r\tau}V^{\rm rec}$ has slope $-rV^{\rm rec}e^{-r\tau}$. On $(T_{\rm par},T_{\rm emb})$ its slope is $e^{-r\tau}(W-rV^{\rm rec})$. At $T_{\rm emb}$ it jumps by $e^{-rT_{\rm emb}}(W/r-V^{\rm rec})$. Under the stated conditions $V^{\rm pay}$ is therefore nondecreasing, and not constant if $rV^{\rm rec}<W$. A higher $\nu_D$ makes $\tau_D$ smaller in the sense of first-order stochastic dominance, so $\mathbb E[V^{\rm pay}(\tau_D)]$ falls, strictly if $V^{\rm pay}$ is not constant. Finally, $\nu_D$ is increasing in $\nu^{\rm L}$ and $\nu^{\rm V}$ and decreasing in $\kappa_{\rm keep}$. For the first-order loss, $\mathbb E[V^{\rm pay}(\tau_D)]=V^{\rm pay}(\infty)e^{-\nu_DT_{\rm emb}}+\int_0^{T_{\rm emb}}\nu_De^{-\nu_D\tau}V^{\rm pay}(\tau)d\tau$, whose derivative at $\nu_D=0$ is $-T_{\rm emb}V^{\rm pay}(\infty)+\int_0^{T_{\rm emb}}V^{\rm pay}(\tau)d\tau$.
(b) $\Pr(\tau_D>T_{\rm emb})=e^{-\nu_DT_{\rm emb}}$. The function $f(T)=\bar vTe^{-r_DT}$ has $f'(T)=\bar ve^{-r_DT}(1-r_DT)$, so it peaks at $1/r_D$. With $r=\nu^{\rm V}=0$, $r_D=\nu^{\rm L}(1-\kappa_{\rm keep})$.
(c) Process $k$ in the order secures $V_{{\rm emb},(k)}$ at its completion date $C_{(k)}$ if no disruption has occurred by then, so the order is worth $\sum_kV_{{\rm emb},(k)}e^{-r_DC_{(k)}}$. Swapping adjacent processes $i$ and $j$ leaves every other term unchanged. Putting $i$ first is at least as good if and only if $V_ie^{-r_D\tau_i}+V_je^{-r_D(\tau_i+\tau_j)}\ge V_je^{-r_D\tau_j}+V_ie^{-r_D(\tau_i+\tau_j)}$, that is, $V_ie^{-r_D\tau_i}(1-e^{-r_D\tau_j})\ge V_je^{-r_D\tau_j}(1-e^{-r_D\tau_i})$, which rearranges to $V_i/(e^{r_D\tau_i}-1)\ge V_j/(e^{r_D\tau_j}-1)$. Any order with an adjacent inversion can be improved weakly by a swap, so the sorted order is optimal. For the switch, let $g(\rho)=(e^{\rho\tau_l}-1)/(e^{\rho\tau_q}-1)$. Then $d\ln g/d\rho=[\phi_e(\rho\tau_l)-\phi_e(\rho\tau_q)]/\rho$ with $\phi_e(u)=u/(1-e^{-u})$, which is strictly increasing, so $g$ is strictly increasing. It rises from $\tau_l/\tau_q$ to $\infty$. The quick win goes first if and only if $V_qg(r_D)>V_l$. Since $V_q\tau_l<V_l\tau_q$, this holds exactly above a unique $r_D^*$.
(d) Differentiate \eqref{eq:ext-Vshock} with respect to $V^{\rm rec}$, and use $\partial V^{\rm rec}/\partial s^{\rm reuse}=F+I$. Both $\nu_D/(r+\nu_D)$ and $1-e^{-(r+\nu_D)T_{\rm emb}}$ increase with $\nu_D$. \qed

\paragraph{Proposition~\ref{prop:ext-continuity}.} (a) Starting from $\mathfrak m=1$ in the distant past,
\[
1-\mathfrak m_t=\sum_{T_i\le t}u_K\,e^{-\eta_K(t-T_i)},\qquad u_K:=\frac{1-\kappa_{\rm doc}}{n_K},
\]
with arrival times $T_i$ at rate $n_K\nu^{\rm K}$. By Campbell's theorem, its mean and variance are
\[
n_K\nu^{\rm K}\int_0^\infty u_Ke^{-\eta_Kt}dt=\frac{\nu^{\rm K}(1-\kappa_{\rm doc})}{\eta_K},\qquad n_K\nu^{\rm K}\int_0^\infty u_K^2e^{-2\eta_Kt}dt=\frac{\nu^{\rm K}(1-\kappa_{\rm doc})^2}{2\eta_Kn_K}.
\] The expected cost flow $L^{\rm ctx}\,\mathbb E[1-\mathfrak m]$ has present value $L^{\rm ctx}\nu^{\rm K}(1-\kappa_{\rm doc})/(\eta_Kr)$, which is linear in $\kappa_{\rm doc}$.
(b) Values are at the overhaul date. Without a device, an informative overhaul replaces the continuation value $-V^-$ by $V^{\rm rec}$, a gain of $V^-+V^{\rm rec}$. An arbitrary one replaces $V^+$ by $V^{\rm rec}$, a loss of $V^+-V^{\rm rec}$. A binding device forgoes both. A gated device keeps the gain and avoids the loss, and its value exceeds the binding device's by $P_Vp_{\rm inf}(V^-+V^{\rm rec})\ge0$.
(c) Up periods last ${\rm Exp}(\nu^{\rm R})$ with mean $1/\nu^{\rm R}$, and freezes last $T_{\rm frz}$. By the renewal-reward theorem the long-run share of time available is $(1/\nu^{\rm R})/(1/\nu^{\rm R}+T_{\rm frz})$.
(d) A wave lasts $T_{\rm par}$ and is hit by a freeze with probability $\nu^{\rm R}T_{\rm par}+O((\nu^{\rm R})^2)$. A freeze in wave $k$ adds $T_{\rm frz}C^{\rm dr}(n_k)$ of dual running. It delays by $T_{\rm frz}$ the savings of every process in waves $k,\dots,n_w$, which costs $T_{\rm frz}W\sum_{i\ge k}n_i$, and it costs $O^{\rm re}_w$ to re-mobilize. Summing over $k$, and using $\sum_k\sum_{i\ge k}n_i=\sum_iin_i$, gives $\nu^{\rm R}T_{\rm frz}T_{\rm par}[\sum_kC^{\rm dr}(n_k)+W\sum_kkn_k]+\nu^{\rm R}T_{\rm par}n_wO^{\rm re}_w$, and the bracket times $T_{\rm par}$ is $\mathcal C^{\rm prog}-n_wO_w$. With equal waves, the cost is
\[
(1+\nu^{\rm R}T_{\rm frz})T_{\rm par}\Bigl[\frac{WN^2}{2n}+\frac{WN}2+c^{\rm dr}_1N+c^{\rm dr}_2Nn\Bigr]+\frac Nn\bigl(O_w+\nu^{\rm R}T_{\rm par}O^{\rm re}_w\bigr),
\]
which is strictly convex in $n$, with the stated minimizer. The change in $n^{*2}$ from $\nu^{\rm R}=0$ has the sign of $(O_w+\nu^{\rm R}T_{\rm par}O^{\rm re}_w)/(1+\nu^{\rm R}T_{\rm frz})-O_w=\nu^{\rm R}(T_{\rm par}O^{\rm re}_w-T_{\rm frz}O_w)/(1+\nu^{\rm R}T_{\rm frz})$. With unit-specific freezes at rate $\nu^{\rm R}/n_s$ per unit, the first-order cost is $\sum_u(\nu^{\rm R}/n_s)\times(\text{stream }u\text{'s cost})$. That is $1/n_s$ of the cost when every freeze, at total rate $\nu^{\rm R}$, stops every stream. \qed

\paragraph{Proposition~\ref{prop:ext-breaks}.} (a) Proposition~\ref{prop:ext-team}(c), together with the intermediate value theorem for continuous, strictly decreasing $\mathcal W(t)$. (b) Proposition~\ref{prop:ext-should}(b). (c) A floor above the current share forces the review share up to $\xi^{\min}(c')$, so the quality cost rises by $c^Q_j(\xi^{\min}(c'))-c^Q_j(\xi^*_j)$. Since $\partial c^Q_j/\partial\xi=\hat w_jt^R_j-e_jL_j\ln(1/\xi)\le\hat w_jt^R_j$, the rise is at most $\hat w_jt^R_j(\xi^{\min}(c')-\xi^*_j)$. Mode A's cost rises by $J_j$, and the non-AI costs are unchanged. If U is exempt, the best AI cost becomes $\min\{\tilde c_{\mathrm A}+J_j,\tilde c_{\mathrm U}\}$, so value falls by $\min\{J_j,\tilde c_{\mathrm U}-\tilde c_{\mathrm A}\}$. (d) Apply (c) with the rise equal to $\mathcal B_j$. (e) to (i) restate the cited results. \qed
